\chapter{工具变量估计}

本书迄今所述的因果推断方法，都依赖一个关键的、无法检验的假设：校正混杂与选择偏倚所需的全部变量都已被识别并正确测量。若这一假设不成立——而它在某种程度上总是不成立——我们的因果估计中就会有残余偏倚。

所幸，另有一类方法可以在另一组假设下有效地估计因果效应，而这组假设不要求测量全部校正因子。工具变量估计就是其中之一。读这本书的经济学家及其他社会科学家现在可以松口气了。我们终于要介绍他们领域中一种极为常见的方法，这种方法与此前讨论过的任何方法都不同。

\section{三个工具变量条件}

图 16.1 中的因果图描绘了一项随机化试验：$Z$ 是随机化分组指示变量（1：治疗组，0：安慰剂组），$A$ 是接受治疗的指示变量（1：是，0：否），因为并非所有参与者都依从自己的分组；$Y$ 是结局；$U$ 是影响结局与依从性的全部因素（其中一些未被测量）。由于参与者及其医生都不知道所服用的药片是治疗还是安慰剂，参与者被称为“盲的”，该研究被称为\emph{双盲安慰剂对照随机化试验}。

\begin{figure}[!ht]
\centering
\begin{tikzpicture}[causal]
  \node[vtx] (Z) at (0,0) {$Z$};
  \node[vtx] (A) at (1.9,0) {$A$};
  \node[vtx] (Y) at (3.8,0) {$Y$};
  \node[vtx] (U) at (0.5,-1.5) {$U$};
  \draw[ed] (Z) -- (A);
  \draw[ed] (A) -- (Y);
  \draw[ed] (U) -- (A);
  \draw[ed] (U) -- (Y);
\end{tikzpicture}
\caption{}
\label{fig:16.1}
\end{figure}

假设我们希望一致地估计 $A$ 对 $Y$ 的平均因果效应。无论使用 IP 加权、标准化、g 估计、分层还是匹配，我们都需要正确测量并校正阻断后门路径 $A\leftarrow U\rightarrow Y$ 的变量，也就是要确保受治疗者与未受治疗者的条件可交换性。不幸的是，如果某些必要变量未被测量、测量不完善，或模型设定有误，上述所有方法都会给出有偏的效应估计。

工具变量（instrumental variable，IV）方法则不同：即便我们没有测量通常所需的、用于校正 $U$ 所致混杂的变量，也可以用它们尝试识别这项随机化试验中 $A$ 对 $Y$ 的平均因果效应。为施展这一魔法，IV 方法需要一个工具变量 $Z$，或称\emph{工具}（instrument）。变量 $Z$ 之所以是工具，是因为它满足三个工具变量条件：

\begin{enumerate}[label=(\roman*),leftmargin=2.4em,itemsep=0.15em,topsep=0.3em]
\item $Z$ 与 $A$ 相关；
\item $Z$ 除了通过其对 $A$ 的潜在效应之外，不影响 $Y$；
\item $Z$ 与 $Y$ 不共享原因。
\end{enumerate}

这些条件的更严格定义见技术注记 16.1，后续的技术注记将沿用该定义。

在上述双盲随机化试验中，随机化指示变量 $Z$ 就是一个工具。条件 (i) 成立，因为被分到治疗组的试验参与者更可能接受治疗；条件 (ii) 由双盲设计所保证；条件 (iii) 由 $Z$ 的随机分配所保证。\mnote{条件 (ii) 得不到保证——例如，参与者因治疗的副作用而在无意中被揭盲。}

\begin{techpoint}[工具变量条件的正式表述]
工具变量条件 (i) 有时称为\emph{相关性条件}（relevance condition），指 $Z$ 与 $A$ 之间存在非零关联，即 $Z\nindep A$ 不成立。在随机化实验中，条件 (i) 预期成立，因为处理分组预期会影响实际接受的治疗。

工具变量条件 (ii) 通常称为\emph{排除限制}（exclusion restriction），指“$Z$ 对 $Y$ 没有直接效应”。在个体层面，条件 (ii) 是：对所有 $z$、$z'$、所有 $a$、所有个体 $i$，都有 $Y_i^{z,a}=Y_i^{z',a}=Y_i^{a}$。然而，本章给出的一些结果只需要总体层面的条件 (ii)，即 $\E[Y^{z,a}]=\E[Y^{z',a}]$。在双盲随机化实验中，两种版本的条件 (ii) 都预期成立，因为若所有个体都不知道分组，分组预期不会影响结局（例如通过行为改变）。对替代工具而言，条件 (ii) 平凡成立。

工具变量条件 (iii) 可以写作\emph{边缘可交换性}：对所有 $a$、$z$，$Y^{a,z}\indep Z$；它在图 16.1、16.2、16.3 对应的 SWIG 中成立。与个体层面的条件 (ii) 合起来，它蕴含 $Y^{a}\indep Z$。更强的一个条件 (iii) 是\emph{联合可交换性}，即对二值处理与二值工具，$\{Y^{z,a};a\in[0,1],z\in[0,1]\}\indep Z$。关于不同类型的可交换性的讨论见技术注记 2.1，关于每一版本可交换性所支撑的结果见技术注记 16.2。在随机化实验中，两种版本的条件 (iii) 都预期成立，因为 $Z$ 是随机分配的。
\end{techpoint}

图 16.1 描绘的是工具 $Z$ 对处理 $A$ 有因果效应的特例。此时我们称 $Z$ 为\emph{因果工具}（causal instrument）。有时因果工具未被测量，我们改用与未测量的因果工具 $U_Z$ 相关的、可测量的\emph{替代工具}（surrogate instrument）$Z$。替代工具对处理 $A$ 没有因果效应，但满足工具变量条件：此时 $Z$ 与 $A$ 的关联 (i) 来自 $Z$ 与 $A$ 的共同原因 $U_Z$，而条件 (iii) 改为“$Z$ 与 $Y$ 除 $U_Z$ 之外不共享原因”。因果工具与替代工具都可用于 IV 估计，但有一些限制，见 16.4 节。作为奇例，图 16.3 描绘了在特定人群中一种不寻常的替代工具 $Z$：$Z$ 与 $A$ 的关联来自对未测量的因果工具 $U_Z$ 与替代工具 $Z$ 的共同效应 $S$ 取条件。

\begin{figure}[!ht]
\centering
\begin{tikzpicture}[causal]
  \node[vtx] (UZ) at (0,0.5) {$U_Z$};
  \node[vtx] (Z) at (0.95,1.65) {$Z$};
  \node[vtx] (A) at (2.45,0.5) {$A$};
  \node[vtx] (Y) at (4.2,0.5) {$Y$};
  \node[vtx] (U) at (0.45,-0.95) {$U$};
  \draw[ed] (UZ) -- (Z);
  \draw[ed] (UZ) -- (A);
  \draw[ed] (Z) -- (A);
  \draw[ed] (A) -- (Y);
  \draw[ed] (U) -- (A);
  \draw[ed] (U) -- (Y);
\end{tikzpicture}
\caption{}
\label{fig:16.2}
\end{figure}

\begin{figure}[!ht]
\centering
\begin{tikzpicture}[causal]
  \node[vtx] (Z) at (0,1.7) {$Z$};
  \node[draw,rectangle,inner sep=1.5pt,minimum size=6mm,font=\small] (S) at (1.5,1.0) {$S$};
  \node[vtx] (UZ) at (0,0.3) {$U_Z$};
  \node[vtx] (A) at (2.6,0.3) {$A$};
  \node[vtx] (Y) at (4.4,0.3) {$Y$};
  \node[vtx] (U) at (0.45,-1.0) {$U$};
  \draw[ed] (Z) -- (S);
  \draw[ed] (UZ) -- (S);
  \draw[ed] (UZ) -- (A);
  \draw[ed] (A) -- (Y);
  \draw[ed] (U) -- (A);
  \draw[ed] (U) -- (Y);
\end{tikzpicture}
\caption{}
\label{fig:16.3}
\end{figure}

此前各章我们使用多种因果推断方法分析观察性数据，估计戒烟对体重变化的影响。要用 IV 方法估计这一效应，我们需要一个工具 $Z$。观察性研究中没有随机化指示变量，因此考虑以下工具候选：香烟价格。可以论证，若 (i) 香烟价格影响戒烟决定，(ii) 香烟价格仅通过其对戒烟的影响而影响体重变化，(iii) 香烟价格与体重变化不存在共同原因，则该变量满足三个工具变量条件。精要 16.1 回顾了观察性研究中提出的一些工具。\mnote{若候选工具 $Z$“出生州的价格”通过未测量的变量 $U_Z$“居住地的价格”与戒烟 $A$ 相关，则条件 (i) 成立。}

\begin{finepoint}[观察性研究中的候选工具变量]
许多变量曾被提出作为观察性研究中的工具，这里无法一一回顾。常用的候选工具分三类：

\begin{itemize}[leftmargin=1.4em,itemsep=0.35em,topsep=0.3em]
\item \textbf{遗传因素}：所提出的工具是一个遗传变异 $Z$，它与处理 $A$ 相关，且据称只通过 $A$ 与结局 $Y$ 有关。例如，在估计饮酒对冠心病风险的影响时，$Z$ 可以是与酒精代谢相关的多态性（比如亚洲人群中的 ALDH2）。利用遗传变异通过 IV 估计从观察性数据作因果推断，属于\emph{孟德尔随机化}（Mendelian randomization）框架（Katan 1986；Davey Smith 与 Ebrahim 2004；Didelez 与 Sheehan 2007；VanderWeele 等 2014）。

\item \textbf{偏好}：所提出的工具 $Z$ 是医生（或医疗服务提供者）对某种治疗相对另一种治疗之偏好的度量。其思路是，医生偏好会影响处方治疗 $A$，而不直接影响结局 $Y$。例如，在估计处方 COX-2 选择性而非非选择性非甾体抗炎药对胃肠道出血的影响时，$U_Z$ 可以是医生对药物类别（COX-2 选择性或非选择性）的处方偏好。由于 $U_Z$ 未被测量，研究者在分析中用（可测量的）替代工具 $Z$ 替代它，例如“医生在当前处方之前开出的上一次处方”（Korn 与 Baumrind 1998；Earle 等 2001；Brookhart 与 Schneeweiss 2007）。

\item \textbf{可及性}：所提出的工具 $Z$ 是治疗可及性的度量。其思路是，可及性影响治疗 $A$ 的使用，但不直接影响结局 $Y$。例如，到某医疗机构的物理距离或出行时间曾被提出作为该机构所提供治疗的工具（McClellan 等 1994；Card 1995；Baiocchi 等 2010）。另一个例子：日历时期曾被提出作为可及性随时间变化的治疗的工具（Hoover 等 1994；Detels 等 1998）。正文中我们用“治疗的价格”——另一种可及性度量——作为候选工具。
\end{itemize}
\end{finepoint}

为使讨论具体，设工具 $Z$ 在个体出生所在美国州的一包香烟平均价格高于 1.50 美元时取 1，否则取 0。遗憾的是，我们无法断定变量 $Z$ 是否真是工具。三个工具变量条件中，只有条件 (i) 可以凭经验检验。要检验它，需要确认所提出的工具 $Z$ 与处理 $A$ 相关，即 $\Pr[A=1|Z=1]-\Pr[A=1|Z=0]>0$。$Z=1$ 者中戒烟的概率为 25.8\%，$Z=0$ 者为 19.5\%；因此风险差 $\Pr[A=1|Z=1]-\Pr[A=1|Z=0]$ 为 6\%。当 $Z$ 与 $A$ 的关联如此微弱时，$Z$ 常被称为\emph{弱工具}（weak instrument）（弱工具的更多讨论见 16.5 节）。

另一方面，条件 (ii) 与 (iii) 无法凭经验检验。要检验条件 (ii)，我们需要证明 $Z$ 只能通过处理 $A$ 引起结局 $Y$。我们不能通过对 $A$ 取条件来证明，因为 $A$ 在图 16.2 的路径 $Z\leftarrow U_Z\rightarrow A\leftarrow U\rightarrow Y$ 上是对撞节点；即使条件 (ii) 为真，对 $A$ 取条件也会在 $Z$ 与 $Y$ 之间诱导出关联。当然，我们也无法证明条件 (iii) 成立，因为我们永远无法排除任何变量之效应的混杂。我们只能假设条件 (ii) 与 (iii) 成立。IV 估计与我们此前研究的所有方法一样，建立在无法检验的假设之上。\mnote{条件 (ii) 与 (iii) 有时可以用工具、处理与结局的数据加以经验否证。然而，否证检验只能排除一小部分违背情形。对大多数违背，该检验没有统计功效，即便样本量任意大（Balke 与 Pearl 1997；Bonet 2001；Glymour 等 2012）。}

在观察性研究中，我们无法证明所提出的工具 $Z$ 真是工具。我们称 $Z$ 为提议的或候选的工具，因为我们永远无法保证图 16.1 与 16.2 所表示的结构就是实际发生的结构。我们能做的最好的事，是运用专业知识为“可以合理地假定所提出的工具 $Z$ 满足条件 (ii) 与 (iii)”建立论证；这与此前各章中我们运用专业知识为各方法的识别假设作辩护是类似的。

但让我们暂且假定 $Z$ 是工具。然后呢？现在我们能见识 IV 估计的魔法如何施展吗？我们能否在不识别、不测量混杂因子的情况下，一致地估计 $A$ 对 $Y$ 的平均因果效应？遗憾的是，若不附加进一步假设，答案是否定的。仅凭一个工具，并不能识别戒烟 $A$ 对体重变化 $Y$ 的平均因果效应，只能识别某些上下界。这些界通常很宽，且往往包含零值（见技术注记 16.2）。

\begin{techpoint}[因果效应的部分识别]
\mnote{界：}对二值结局 $Y$，平均因果效应可在 $-1$（若所有人除非受治疗否则都不发生该结局）与 $1$（若无人除非受治疗否则发生该结局）之间取值。平均因果效应的界为 $(-1,1)$。利用数据可以把这两个界之间的距离缩小一半：因为对每个个体，我们知道其两个反事实结局之一——若该个体实际受治疗，则知道 $Y^{a=1}$；若实际未受治疗，则知道 $Y^{a=0}$——所以我们可以对每个个体未知的反事实结局赋予最极端的可能取值，再计算因果效应。这样得到的平均因果效应的界更窄，但仍包含零值 0。对连续结局 $Y$，推导界需要指定结局的最小值与最大值；界的宽度随所选取值而变。

$\Pr[Y^{a=1}=1]-\Pr[Y^{a=0}=1]$ 的界可以进一步收窄，只要存在变量 $Z$ 满足总体层面的工具变量条件 (ii)（见技术注记 16.1）与边缘可交换性 (iii)（Robins 1989；Manski 1990）。这些所谓\emph{自然界}的宽度为 $\Pr[A=1|Z=0]+\Pr[A=0|Z=1]$，比仅凭数据识别出的界更窄。有时，若把边缘可交换性换成联合可交换性，还可以得到更窄的界——所谓\emph{锐界}（sharp bounds）（Balke 与 Pearl 1997；Richardson 与 Robins 2014）。

达到锐界所需的条件，也可以从对 $z$ 与 $a$ 作联合干预的 SWIG 推出，它们对应于图 16.1、16.2、16.3 中任一因果图。Richardson 与 Robins (2010, 2014) 证明：条件 $Y^{a,z}\indep(Z,A)\given U$ 与 $Z\indep U$，再加上 $U$ 各层内的总体层面条件 (ii)，即 $\E[Y^{z,a}|U]=\E[Y^{z',a}|U]$，足以得到锐界。这些条件对三个 SWIG 都成立，它们蕴含 $Z\indep U$ 与 $Y\indep Z\given U,A$，并蕴含 $\E[Y^{z,a}]$ 由忽略 $Z$ 的 g 公式 $\int \E[Y|A=a,U=u]\,dF(u)$ 给出；这反映了在 $U$ 各层内 $Z$ 对 $Y$ 没有直接效应。Dawid (2003) 证明后面这些条件可导出锐界。在进一步假设下，Richardson 与 Robins 导出了更窄的界。另见 Richardson、Evans 与 Robins (2011)。

遗憾的是，所有这些部分识别方法（即界定效应的方法）往往信息量不大，因为界很宽。Swanson 等 (2018) 综述了针对二值工具、处理与结局的部分识别方法。Swanson 等 (2015a) 描述了几种部分识别方法在真实数据中的应用，并讨论了各自的优缺点。

有一种办法可以减小界的宽度：对 $A$ 对 $Y$ 之效应的形式作参数假设。在 16.2 节所述的足够强的假设下，上下界收敛为一个数，平均因果效应被点识别。
\end{techpoint}

在我们的例子中，这些界帮助不大。它们只会确认我们已知的事：戒烟可能导致体重增加、体重减少，或体重不变。遗憾的是，除非愿意作出额外的、无法验证的假设，工具能提供的只有这样。16.3 节与 16.4 节回顾使 IV 估计目标成为平均因果效应的附加条件。在此之前，我们先回顾具体做法。

\section{通常的工具变量估计目标}

当一个二值变量 $Z$ 是工具，即满足三个工具变量条件 (i)–(iii)，并且下一节所述的附加条件 (iv) 也成立时，处理对结局在加法尺度上的平均因果效应 $\E[Y^{a=1}]-\E[Y^{a=0}]$ 可被识别，且等于
\[
\frac{\E[Y|Z=1]-\E[Y|Z=0]}{\E[A|Z=1]-\E[A|Z=0]},
\]
这就是二值工具下\emph{通常的 IV 估计目标}。（对二值处理，注意 $\E[A|Z=1]=\Pr[A=1|Z=1]$。）技术注记 16.3 用可加结构均值模型给出这一结果的证明，但你或许愿意先读下一节，再回头看它。\mnote{我们将集中讨论二值工具，这是最常见的情形。对连续工具 $Z$，通常的 IV 估计目标为 $\mathrm{Cov}(Y,Z)$ 与 $\mathrm{Cov}(A,Z)$ 之比，其中 Cov 表示协方差。}

为直观理解通常的 IV 估计目标，再考虑上一节的随机化试验。IV 估计目标的分子是分组 $Z$ 对 $Y$ 的平均因果效应——即\emph{意向性治疗效应}；分母是分组 $Z$ 对 $A$ 的平均因果效应——即对分派治疗的依从程度（adherence，或称 compliance）的度量。若依从完美，分母等于 1，$A$ 对 $Y$ 的效应就等于 $Z$ 对 $Y$ 的效应。依从性越差，分母越接近 0，$A$ 对 $Y$ 的效应就越大。依从性越低，$A$ 对 $Y$ 的效应——即 IV 估计目标——与 $Z$ 对 $Y$ 的效应之间的差距就越大。\mnote{在随机化实验中，IV 估计目标是两个效应之比：$Z$ 对 $Y$ 的效应与 $Z$ 对 $A$ 的效应。由于 $Z$ 是随机分配的，这两个效应都可以不做校正而一致地估计。}

IV 估计目标通过放大分组的效应（分子），绕过了校正混杂因子的需要。放大的幅度随依从性下降而增大，即随 $Z$ 与 $A$ 的风险差（分母）趋近于零而增大。同一道理适用于观察性研究中所用的工具：此时 IV 估计量的分母可能等于因果工具 $Z$ 对 $A$ 的因果效应（图 16.1），也可能等于替代工具 $Z$ 与处理 $A$ 之间的非因果关联（图 16.2 与 16.3）。

标准的 IV 估计量，按通常的 IV 估计目标之分子与分母的估计值之比计算。\mnote{亦称 Wald 估计量（Wald 1940）。}在戒烟例子中，工具 $Z$ 为二值（1：高香烟价格的州，0：其他），分子估计 $\hat{\E}[Y|Z=1]-\hat{\E}[Y|Z=0]$ 等于 $2.686-2.536=0.1503$，分母估计 $\hat{\E}[A|Z=1]-\hat{\E}[A|Z=0]$ 等于 $0.2578-0.1951=0.0627$。因此通常的 IV 估计值为比值 $0.1503/0.0627=2.4$ kg。在三个工具变量条件 (i)–(iii) 加上下一节的条件 (iv) 之下，这就是总体中戒烟对体重增加之平均因果效应的估计。

我们只是计算了四个样本均值 $\hat{\E}[A|Z=1]$、$\hat{\E}[A|Z=0]$、$\hat{\E}[Y|Z=1]$、$\hat{\E}[Y|Z=0]$，就得到了 IV 估计目标之分子与分母的估计。等价地，我们也可以拟合两个（饱和）线性模型，来估计分母与分子中的差值。分母的模型为 $\E[A|Z]=\alpha_0+\alpha_1 Z$，分子的模型为 $\E[Y|Z]=\beta_0+\beta_1 Z$。\mnote{代码：程序 16.1。这里排除结局或治疗缺失者；实践中可先用 IP 加权校正选择偏倚。}

计算标准 IV 估计量的另一种方法，是\emph{两阶段最小二乘估计量}（two-stage-least-squares estimator）。步骤如下。首先，拟合第一阶段治疗模型 $\E[A|Z]=\alpha_0+\alpha_1 Z$，并为每个个体生成预测值 $\hat{\E}[A|Z]$。其次，拟合第二阶段结局模型 $\E[Y|Z]=\beta_0+\beta_1\hat{\E}[A|Z]$。参数估计 $\hat\beta_1$ 在数值上总是等于标准 IV 估计值。因此在本例中，两阶段最小二乘估计同样是 2.4 kg。\mnote{代码：程序 16.2}

2.4 这一点估计的 95\% 置信区间非常大：$-36.5$ 至 $41.3$。对我们所提出的工具而言，这是意料之中的，因为 $Z$ 与 $A$ 的关联很弱，第一阶段模型的不确定性很大。一条常用的经验法则是：若第一阶段模型的 $F$ 统计量小于 10，就宣布工具为弱工具（本例中它仅为 0.8）。我们将在 16.5 节再讨论弱工具带来的问题。

把两阶段最小二乘估计量视为识别处理之因果效应，其中隐含的一些假设，可以通过使用可加或可乘结构均值模型（如技术注记 16.3 与 16.4 所述）作 IV 估计而变得更明确。结构均值模型的参数可用 g 估计来估计。此外，当存在可测量的、工具与结局的共同原因 $L$、因而分析中必须加以校正时，在两阶段最小二乘线性模型与结构均值模型之间作取舍所涉及的权衡，可能与在非 IV 估计中于结局回归与结构嵌套模型之间作取舍所涉及的权衡类似（见第 14、15 章）。\mnote{代码：程序 16.3}

无论如何，上述估计量只有在通常的 IV 估计目标可以解释为处理 $A$ 对结局 $Y$ 的平均因果效应时才有效。为此，除三个工具变量条件外，还需要第四个识别条件成立。

\section{第四个识别条件：同质性}

三个工具变量条件 (i)–(iii) 不足以保证 IV 估计目标就是处理 $A$ 对 $Y$ 的平均因果效应。还需要第四个条件——\emph{效应同质性}（effect homogeneity）(iv)。下面按（历史上的）出现顺序描述四种可能的同质性条件 (iv)。

最极端、也最古老的一种同质性条件 (iv)，是处理 $A$ 对结局 $Y$ 的效应在个体间恒定。在我们的例子中，若戒烟使总体中每个人体重增加（或减少）同样的量，比如恰好 2.4 kg，该条件就成立。恒定效应等价于\emph{可加秩保持}（additive rank preservation）；正如 14.4 节所论，对多数处理与结局而言，这在科学上不可信——对二值结局则不可能，除非在锐零假设或普遍受害（或受益）的情形下。在我们的例子中，我们预计戒烟后有些人会大幅增重，有些人增重很少，还有人甚至会减重。因此，我们一般不愿接受恒定效应这一同质性假设作为合理的条件 (iv)。\mnote{然而，在许多早期使用两阶段最小二乘估计量的 IV 分析中，可加秩保持被隐含地假定了。}

对二值 $Z$ 与 $A$，第二种不那么极端的同质性条件 (iv) 是：$A$ 对 $Y$ 的平均因果效应在 $Z$ 的各层之间相等，且在接受治疗者与未接受治疗者中都如此，即
\[
\E[Y^{a=1}-Y^{a=0}|Z=1,A=a]=\E[Y^{a=1}-Y^{a=0}|Z=0,A=a],\qquad a=0,1.
\]
这一\emph{可加}同质性条件 (iv) 正是技术注记 16.3 的数学证明中所用的条件。乘法尺度上的另一种同质性条件见技术注记 16.4。（该可乘同质性条件导出的 IV 估计目标与通常的 IV 估计目标不同。）

上述同质性条件的表述并不自然直观。\mnote{即便 $Y^{a}\indep Z$ 成立（如图 16.1–16.3 的 SWIG），$Y^{a=1}-Y^{a=0}\indep Z$ 一般也不成立，故效应可能依赖 $Z$。}在具体研究中，专家如何为“在所提出的工具 $Z$ 与处理 $A$ 的各层内平均因果效应恒定”提供论据？更自然的同质性条件 (iv)——尽管仍然无法检验——会以可能已知（即便未被测量）的混杂因子 $U$ 所致效应修饰来表述。一种这样的条件是：$U$ 不是可加效应修饰因子，即 $A$ 对 $Y$ 的平均因果效应在未测量混杂因子 $U$ 的每一层都相同，$\E[Y^{a=1}|U]-\E[Y^{a=0}|U]=\E[Y^{a=1}]-\E[Y^{a=0}]$。这第三种同质性条件 (iv) 往往不可信，因为某些未测量的混杂因子可能同时也是效应修饰因子。例如，戒烟后体重增加的幅度可能随既往吸烟强度而异，而既往吸烟强度本身可能是戒烟对体重增加之效应的未测量混杂因子。\mnote{Hernán 与 Robins (2006b) 表明，若 $U$ 是效应修饰因子，那么我们就不太可能有理由相信前述可加同质性条件 (iv) 成立。}

另一类同质性条件 (iv) 是：$Z$ 与 $A$ 的关联在加法尺度上跨未测量混杂因子 $U$ 的各层恒定，即 $\E[A|Z=1,U]-\E[A|Z=0,U]=\E[A|Z=1]-\E[A|Z=0]$。这一条件与前述条件 $\E[Y^{a=1}|U]-\E[Y^{a=0}|U]=\E[Y^{a=1}]-\E[Y^{a=0}]$，都是下述（远为）一般的条件的特例：$U$ 对处理 $A$ 影响结局 $Y$ 之效应的修饰 $e(U)\equiv\E[Y^{a=1}-Y^{a=0}|U]$，与 $U$ 对 $Z$ 与 $A$ 之关联在加法尺度上的修饰 $t(U)\equiv\E[A|Z=1,U]-\E[A|Z=0,U]$ 不相关，即 $\mathrm{Cov}[e(U),t(U)]=0$。前两个条件都是 $\mathrm{Cov}[e(U),t(U)]=0$ 的特例，因为它们分别可表述为“$e(U)$ 是常数”与“$t(U)$ 是常数”，而任何变量与常数的协方差总为 0。\mnote{同质性条件“$t(U)$ 是常数”是一般条件的一个特例（Wang 与 Tchetgen Tchetgen 2018；Hartwig 等 2023）。证明见技术注记 16.5。}

\begin{techpoint}[可加结构均值模型与 IV 估计]
考虑如下针对二值处理 $A$ 与工具 $Z$（如图 16.1、16.2 或 16.3 所示）的饱和可加结构均值模型：
\[
\E[Y^{a=1}-Y^{a=0}|A=1,Z]=\beta_0+\beta_1 Z.
\]
该模型也可写作 $\E[Y-Y^{a=0}|A,Z]=A(\beta_0+\beta_1 Z)$。参数 $\beta_0$ 是 $Z=0$ 的受治疗者中处理的平均因果效应，$\beta_0+\beta_1$ 是 $Z=1$ 的受治疗者中处理的平均因果效应。因此 $\beta_1$ 量化了 $Z$ 的可加效应修饰。

若我们先验地假定不存在 $Z$ 的可加效应修饰，则 $\beta_1=0$，$\beta_0$ 恰好就是通常的 IV 估计目标（Robins 1994）。也就是说，通常的 IV 估计目标，是在不存在 $Z$ 所致效应修饰时，处理对受治疗者之效应的可加结构均值模型的参数。

证明很简单。当 $Z$ 是工具时，条件 (ii) 成立，蕴含 $\E[Y^{a=0}|Z=1]=\E[Y^{a=0}|Z=0]$。在上述结构模型下，这一条件均值独立性可改写为 $\E[Y-A(\beta_0+\beta_1)|Z=1]=\E[Y-A\beta_0|Z=0]$。在 $\beta_1=0$ 下解此方程，得
\[
\beta_0=\frac{\E[Y|Z=1]-\E[Y|Z=0]}{\E[A|Z=1]-\E[A|Z=0]}.
\]

你或许会问，我们为什么先验地令 $\beta_1=0$。原因在于，我们有一个含两个未知量（$\beta_0$ 与 $\beta_1$）的方程，而该方程已穷尽了三个工具变量条件对数据分布所施加的全部约束。既然还需要一个额外约束，而按定义它无法检验，我们便任意地选择 $\beta_1=0$（而不是比如 $\beta_1=2$）。这就是我们说“仅凭工具不足以识别平均因果效应”的意思。

要断定受治疗者中处理的平均因果效应 $\beta_0=\E[Y^{a=1}-Y^{a=0}|A=1]$ 等于研究总体中的平均因果效应 $\E[Y^{a=1}]-\E[Y^{a=0}]$——从而断定通常的 IV 估计目标就是 $\E[Y^{a=1}]-\E[Y^{a=0}]$——我们必须假定受治疗者与未受治疗者的处理效应相同，这是一个额外的、无法检验的假设。于是，在附加假设 $\beta_1=0$ 下，对任意 $z$，$\beta_0=\E[Y^{a=1}-Y^{a=0}|A=1,Z=z]=\E[Y^{a=1}-Y^{a=0}|A=1]$ 就是受治疗者中处理的平均因果效应。要断定 $\beta_0$ 是研究总体中的平均因果效应 $\E[Y^{a=1}]-\E[Y^{a=0}]$——从而断定它就是通常的 IV 估计目标——我们还必须假定处理效应在受治疗者与未受治疗者中相同，即 $A=0$ 的结构模型中 $Z$ 的参数也为 0。这是又一个额外的、无法检验的假设。
\end{techpoint}

\begin{techpoint}[可乘结构均值模型与 IV 估计]
考虑如下针对二值处理 $A$ 的饱和可乘（对数线性）结构均值模型：
\[
\frac{\E[Y^{a=1}|A=1,Z]}{\E[Y^{a=0}|A=1,Z]}=\exp(\beta_0+\beta_1 Z),
\]
它也可写作 $\E[Y|A,Z]=\E[Y^{a=0}|A,Z]\exp[A(\beta_0+\beta_1 Z)]$。对二值 $Y$，$\exp(\beta_0)$ 是 $Z=0$ 的受治疗者中的因果风险比，$\exp(\beta_0+\beta_1)$ 是 $Z=1$ 的受治疗者中的因果风险比。因此 $\beta_1$ 量化了 $Z$ 的可乘效应修饰。若我们先验地假定 $\beta_1=0$——并额外假定未受治疗者中也不存在 $Z$ 的可乘效应修饰——则乘法（风险比）尺度上的因果效应为 $\E[Y^{a=1}]/\E[Y^{a=0}]=\exp(\beta_0)$，加法（风险差）尺度上的因果效应为
\begin{align*}
\E[Y^{a=1}]-\E[Y^{a=0}]={}&\E[Y|A=0](1-\E[A])[\exp(\beta_0)-1]\\
&+\E[Y|A=1]\E[A][1-\exp(-\beta_0)].
\end{align*}
证明依赖工具变量条件，见 Robins (1989) 与 Hernán、Robins (2006b)。

也就是说，若我们假定一个在受治疗者与未受治疗者中都不存在 $Z$ 所致可乘效应修饰的可乘结构均值模型，则平均因果效应 $\E[Y^{a=1}]-\E[Y^{a=0}]$ 仍可识别，但不再等于通常的 IV 估计量。因此，我们对 $\E[Y^{a=1}]-\E[Y^{a=0}]$ 的估计将取决于我们假定不存在 $Z$ 的\emph{可加}效应修饰还是\emph{可乘}效应修饰。遗憾的是，即便样本量无穷大，也无法确定这两个假设（若有的话）中哪个为真（Robins 1994），因为在考虑饱和的可加或可乘结构均值模型时，待估的未知参数多于用来估计它们的方程。这正是我们需要作出同质性之类建模假设的原因。
\end{techpoint}

\begin{techpoint}[一般同质性条件的证明]
我们希望证明：在图 16.1 的因果图与零协方差 $\mathrm{Cov}(e(U),t(U))=0$ 的一般同质性条件下，
{\footnotesize
\[
\E[Y^{a=1}-Y^{a=0}]=\frac{\E[Y|Z=1]-\E[Y|Z=0]}{\E[A|Z=1]-\E[A|Z=0]},
\]}
其中 $e(U)$ 与 $t(U)$ 如正文所定义。

为此，注意 $\E[e(U)]=\E[Y^{a=1}-Y^{a=0}]$。又因 $U\indep Z$，$\E[t(U)]=\E[A|Z=1]-\E[A|Z=0]$。因此零协方差条件蕴含 $\E[e(U)t(U)]/\{\E[A|Z=1]-\E[A|Z=0]\}=\E[Y^{a=1}-Y^{a=0}]$。尚需证明 $\E[Y|Z=1]-\E[Y|Z=0]=\E[e(U)t(U)]$。为此，把 $Y$ 写作 $Y=A(Y^{a=1}-Y^{a=0})+Y^{a=0}$。由于图 16.1 中 $Y^{a}\indep(A,Z)\given U$ 且 $U\indep Z$，$\E[Y|Z]$ 等于
\begin{align*}
\sum_{a\in\{0,1\}}\int &\E[Y|A=a,Z,U=u]\Pr(A=a|Z,U=u)\,f(u|Z)\,du\\
&=\int\Bigl\{\E[Y^{a=1}-Y^{a=0}|U=u]\Pr(A=1|Z,U=u)\\
&\qquad+\E[Y^{a=0}|U=u]\Bigr\}f(u)\,du.
\end{align*}
因此
\begin{align*}
\E[Y|Z=1]-\E[Y|Z=0]=\E\Bigl[\E[Y^{a=1}-Y^{a=0}|U]\bigl\{&\Pr(A=1|Z=1,U)\\
&-\Pr(A=1|Z=0,U)\bigr\}\Bigr],
\end{align*}
如所要求。
\end{techpoint}

鉴于同质性条件在许多情形下明显不可信，IV 方法能否有效地估计处理的平均因果效应，似乎值得怀疑。有两种做法可以绕过同质性条件。

一种做法是在 IV 估计的模型中引入基线协变量。为此，使用结构均值模型更稳妥，因为它比两阶段最小二乘估计量施加的参数假设更少。在结构均值模型中纳入协变量，通过对处理效应在协变量各层内如何变化施加约束，允许受治疗者的处理效应随 $Z$ 而变。关于带协变量的结构均值模型的更多细节，见 16.5 节与技术注记 16.6。

另一种做法是采用不需要效应同质性的替代条件 (iv)。与三个工具变量条件 (i)–(iii) 结合时，这一替代条件使我们能给通常的 IV 估计目标赋予因果解释，尽管它不足以识别总体中的平均因果效应。下一节回顾这一替代条件 (iv)。

\section{另一种第四条件：单调性}

再次考虑以随机化指示变量 $Z$、处理 $A$ 与结局 $Y$ 构成的双盲随机化试验。对试验中的每个个体，反事实变量 $A^{z=1}$ 是“若被分派接受治疗（$z=1$）时该个体会采取的处理取值”——1 或 0。反事实变量 $A^{z=0}$ 类似地定义为“若被分派不接受治疗（$z=0$）时的处理取值”。

若我们知道每个个体的两个反事实处理变量 $A^{z=1}$ 与 $A^{z=0}$ 的取值，就可以把研究总体中的全部个体分为四个互不相交的子总体：

\begin{enumerate}[leftmargin=1.9em,itemsep=0.4em,topsep=0.3em]
\item \textbf{永远接受者}（always-takers）：无论被分到哪一组都会接受治疗的个体。即同时有 $A^{z=1}=1$ 与 $A^{z=0}=1$ 的个体。

\item \textbf{永不接受者}（never-takers）：无论被分到哪一组都不接受治疗的个体。即同时有 $A^{z=1}=0$ 与 $A^{z=0}=0$ 的个体。

\item \textbf{依从者}（compliers）或称\emph{合作者}（cooperative）：被分到治疗组就接受治疗、被分到不治疗组就不接受治疗的个体。即 $A^{z=1}=1$ 且 $A^{z=0}=0$ 的个体。

\item \textbf{违抗者}（defiers）或称\emph{对抗者}（contrarians）：被分到治疗组却不接受治疗、被分到不治疗组却接受治疗的个体。即 $A^{z=1}=0$ 且 $A^{z=0}=1$ 的个体。
\end{enumerate}

注意，这些子总体——常被称为\emph{依从类型}（compliance types）或\emph{主层}（principal strata）——一般不可识别。若我们观察到某人被分到 $Z=1$ 且接受了治疗 $A=1$，我们并不知道她是依从者还是永远接受者。若观察到某人被分到 $Z=1$ 且未接受治疗 $A=0$，我们也不知道他是违抗者还是永不接受者。

当不存在违抗者时，我们说\emph{单调性}成立，因为工具 $Z$ 要么不改变处理 $A$——如图 16.4 的永远接受者与图 16.5 的永不接受者——要么提高处理 $A$ 的取值——如图 16.6 的依从者。对违抗者，工具 $Z$ 会降低处理 $A$ 的取值——如图 16.7 所示。更一般地，当对所有个体都有 $A^{z=1}\ge A^{z=0}$ 时，单调性成立。

\begin{figure}[!ht]
\centering
\begin{tikzpicture}[causal]
  \draw[gray] (0,0)--(2.9,0);
  \draw[gray] (0,0)--(0,1.95);
  \draw[gray] (0.9,0.08)--(0.9,-0.08);
  \draw[gray] (2.1,0.08)--(2.1,-0.08);
  \draw[gray] (0.08,0.45)--(-0.08,0.45);
  \draw[gray] (0.08,1.45)--(-0.08,1.45);
  \node[left,font=\footnotesize] at (-0.12,0.45) {$0$};
  \node[left,font=\footnotesize] at (-0.12,1.45) {$1$};
  \node[left,font=\footnotesize] at (-0.55,0.95) {$A^{z}$};
  \node[below,font=\footnotesize] at (0.9,-0.1) {$z=0$};
  \node[below,font=\footnotesize] at (2.1,-0.1) {$z=1$};
  \draw[very thick] (0.9,1.45)--(2.1,1.45);
  \fill (0.9,1.45) circle (1.5pt);
  \fill (2.1,1.45) circle (1.5pt);
  \node[above,font=\footnotesize] at (1.5,1.62) {永远接受者};
\end{tikzpicture}
\caption{}
\label{fig:16.4}
\end{figure}

\begin{figure}[!ht]
\centering
\begin{tikzpicture}[causal]
  \draw[gray] (0,0)--(2.9,0);
  \draw[gray] (0,0)--(0,1.95);
  \draw[gray] (0.9,0.08)--(0.9,-0.08);
  \draw[gray] (2.1,0.08)--(2.1,-0.08);
  \draw[gray] (0.08,0.45)--(-0.08,0.45);
  \draw[gray] (0.08,1.45)--(-0.08,1.45);
  \node[left,font=\footnotesize] at (-0.12,0.45) {$0$};
  \node[left,font=\footnotesize] at (-0.12,1.45) {$1$};
  \node[left,font=\footnotesize] at (-0.55,0.95) {$A^{z}$};
  \node[below,font=\footnotesize] at (0.9,-0.1) {$z=0$};
  \node[below,font=\footnotesize] at (2.1,-0.1) {$z=1$};
  \draw[very thick] (0.9,0.45)--(2.1,0.45);
  \fill (0.9,0.45) circle (1.5pt);
  \fill (2.1,0.45) circle (1.5pt);
  \node[above,font=\footnotesize] at (1.5,0.62) {永不接受者};
\end{tikzpicture}
\caption{}
\label{fig:16.5}
\end{figure}

\begin{figure}[!ht]
\centering
\begin{tikzpicture}[causal]
  \draw[gray] (0,0)--(2.9,0);
  \draw[gray] (0,0)--(0,1.95);
  \draw[gray] (0.9,0.08)--(0.9,-0.08);
  \draw[gray] (2.1,0.08)--(2.1,-0.08);
  \draw[gray] (0.08,0.45)--(-0.08,0.45);
  \draw[gray] (0.08,1.45)--(-0.08,1.45);
  \node[left,font=\footnotesize] at (-0.12,0.45) {$0$};
  \node[left,font=\footnotesize] at (-0.12,1.45) {$1$};
  \node[left,font=\footnotesize] at (-0.55,0.95) {$A^{z}$};
  \node[below,font=\footnotesize] at (0.9,-0.1) {$z=0$};
  \node[below,font=\footnotesize] at (2.1,-0.1) {$z=1$};
  \draw[very thick] (0.9,0.45)--(2.1,1.45);
  \fill (0.9,0.45) circle (1.5pt);
  \fill (2.1,1.45) circle (1.5pt);
  \node[font=\footnotesize] at (1.35,1.2) {依从者};
\end{tikzpicture}
\caption{}
\label{fig:16.6}
\end{figure}

\begin{figure}[!ht]
\centering
\begin{tikzpicture}[causal]
  \draw[gray] (0,0)--(2.9,0);
  \draw[gray] (0,0)--(0,1.95);
  \draw[gray] (0.9,0.08)--(0.9,-0.08);
  \draw[gray] (2.1,0.08)--(2.1,-0.08);
  \draw[gray] (0.08,0.45)--(-0.08,0.45);
  \draw[gray] (0.08,1.45)--(-0.08,1.45);
  \node[left,font=\footnotesize] at (-0.12,0.45) {$0$};
  \node[left,font=\footnotesize] at (-0.12,1.45) {$1$};
  \node[left,font=\footnotesize] at (-0.55,0.95) {$A^{z}$};
  \node[below,font=\footnotesize] at (0.9,-0.1) {$z=0$};
  \node[below,font=\footnotesize] at (2.1,-0.1) {$z=1$};
  \draw[very thick] (0.9,1.45)--(2.1,0.45);
  \fill (0.9,1.45) circle (1.5pt);
  \fill (2.1,0.45) circle (1.5pt);
  \node[font=\footnotesize] at (1.45,1.25) {违抗者};
\end{tikzpicture}
\caption{}
\label{fig:16.7}
\end{figure}

现在我们用单调性条件替换上一节的任何一种同质性条件，它将成为我们的新条件 (iv)。于是通常的 IV 估计目标不再等于处理的平均因果效应 $\E[Y^{a=1}]-\E[Y^{a=0}]$。相反，在单调性 (iv) 下，通常的 IV 估计目标等于依从者中处理的平均因果效应，即
\[
\E[Y^{a=1}-Y^{a=0}|A^{z=1}=1,A^{z=0}=0].
\]
技术注记 16.7 在 $Z$ 被有效随机分配的假设下给出了这一等式的证明。作为证明的梗概：通常的 IV 估计目标与依从者中的效应之所以相等，是因为分组 $Z$ 对 $Y$ 的效应——IV 估计目标的分子——是 $Z$ 在四个主层各自效应的加权平均。然而，$Z$ 对 $Y$ 的效应在永远接受者与永不接受者中恰好为零，因为 $Z$ 的效应完全经由 $A$ 中介，而这些子总体中 $A$ 的取值是固定的，无论被分到 $Z$ 的哪个取值。此外，在单调性 (iv) 下不存在违抗者。因此 IV 估计目标的分子是 $Z$ 对依从者中 $Y$ 的效应——它等同于 $A$ 对依从者中 $Y$ 的效应——乘以总体中依从者的比例，而后者恰是通常的 IV 估计目标的分母。

在观察性研究中，通常的 IV 估计目标也可以用来估计无违抗者时依从者中的效应。严格说来，观察性研究中没有依从者与违抗者，因为所提出的工具 $Z$ 不是处理分派；但“依从者”一词指的是满足 $(A^{z=1}=1,A^{z=0}=0)$ 的个体，“违抗者”指满足 $(A^{z=1}=0,A^{z=0}=1)$ 的个体。在戒烟例子中，依从者是那些在高香烟价格的州会戒烟、在低价格的州不会戒烟的人。反之，违抗者是那些在高香烟价格的州不会戒烟、在低价格的州会戒烟的人。若不存在违抗者且因果工具是二值的（见下文及技术注记 16.6），则 2.4 kg 就是依从者中的 IV 效应估计值。

以单调性取代同质性，在 1990 年代中期被欢迎为 IV 方法的救星。同质性往往不可信，而单调性在许多情形下显得可信。单调性 (iv) 下的 IV 方法不能识别总体中的平均因果效应，只能识别依从者子总体中的效应，但为了把强大的 IV 方法留在工具箱里，这个代价似乎值得。然而，在单调性 (iv) 下估计依从者中处理的平均因果效应，已受到多方面的批评。

第一，依从者中效应的相关性值得怀疑。\marginnote{\color{black!85}“依从者平均因果效应”（CACE）是子总体中的局部平均因果效应（LATE），不是总体的全局效应。Greenland (2000) 称之为“依从者-依从”效应。\endgraf\vspace{2pt}Deaton (2010)：这好比任光线落在何处，再宣称凡被照亮者正是我们一直在找的东西。}依从者子总体不可识别；尽管依从者在总体中的比例可以计算（它就是通常的 IV 估计目标的分母，见技术注记 16.7），但它因工具而异、因研究而异。因此，关于依从者中效应的因果推断难以被决策者使用。既然处理已被估计为对依从者有益——在本例中他们恰是总体的 6\%，但在真实世界中这一群体可能更小或更大——决策者是否就应优先向整个人群实施 $A=1$？如果处理对永远接受者与永不接受者（他们占人口多数）并不那么有益呢？遗憾的是，决策者无法知道谁属于这 6\%。与其争辩依从者的效应是首要关切的，不如承认：人们对这一估计目标的兴趣，并非源于其实用相关性，而是源于一种（往往错误的）看法——它容易识别。

第二，单调性在观察性研究中并不总是合理的假设。在随机化试验中，不存在违抗者似乎是稳妥的假设：我们不预期有些人会同意参加试验，却怀着相反的意图，偏要做与要求相反的事。此外，在那些阻止被分到不治疗组者接受治疗的试验中，单调性由设计所保证，即不存在永远接受者或违抗者。在这种情形下，依从者中的效应实际上就是受治疗者中的效应。\mnote{一个缓解因素是，在较强的假设下，研究者可以用观测变量的分布来刻画依从者（Angrist 与 Pischke 2009；Baiocchi 等 2014）。}

然而，对观察性研究中提出的一些工具，单调性较难辩护。考虑以“医生偏好”作为候选工具，估计在两位偏好不同的医生坐诊的诊所就诊的患者的处理效应。第一位医生通常倾向于开这种治疗，但对患糖尿病的患者例外（因为某些已知的禁忌证）。第二位医生通常倾向于不开这种治疗，但对更常运动的患者例外（因为某些已感知的益处）。任何既常运动又患糖尿病的患者，都会违背这两位医生的偏好而接受治疗，因而会被标为违抗者。也就是说，当治疗决定是权衡多个标准或多个鼓励维度（其中既有风险也有收益）的结果时，单调性不太可能成立。在这些情形下，违抗者的比例可能不可忽略。\mnote{上面这个例子由 Swanson 与 Hernán (2014) 提出。Swanson 等 (2015b) 还在观察性情形中实证地展示了违抗者的存在。}

对图 16.2 与 16.3 所表示的替代工具 $Z$，情况更为复杂。若因果工具 $U_Z$ 是连续的（例如真实的、未测量的医生偏好），则使用二值替代工具 $Z$（例如偏好的某种可测量替代）的标准 IV 估计目标，并不是某一特定依从者子总体中的效应。相反，标准 IV 估计目标识别的是总体中所有个体之效应的某个特定加权平均，这使其难以解释。因此，当所提出的二值工具不是因果工具时，即便单调性对连续的因果工具 $U_Z$ 成立，把 IV 估计目标解释为依从者中的效应也值得怀疑（细节见技术注记 16.7）。\mnote{连续因果工具 $U_Z$ 的单调性定义：$A^{u_Z}$ 是 $u_Z$ 在 $U_Z$ 支撑集上的非减函数（Angrist 与 Imbens 1995；Heckman 与 Vytlacil 1999）。}

最后但绝非最不重要的一点，把总体划分为四个子总体或主层，可能并不正当。在许多现实情形中，依从者子总体是一个定义不良的总体子集。例如，在有多位医生的情形中使用“医生偏好”作为候选工具，所有偏好水平相同、且本可能接诊某位患者的医生，都必须以完全相同的方式治疗该患者。这不仅是不切实际的假设，而且在许多观察性研究中根本不可能定义，因为我们不知道哪些医生本可能接诊该患者。要把总体稳定地划分为依从者、违抗者、永远接受者与永不接受者，还需要确定性反事实（估计平均因果效应一般并不要求这一点）、无干扰（例如，我可能是永远接受者，但若我的朋友不接受治疗，我也会决定不接受），以及若干形式的异质性（在某种情形下、或对某个工具而言的依从者，在另一种情形下可能不是依从者）。\mnote{Swanson 等 (2015b) 讨论了定义单调性的困难，并引入了观察性研究中全局单调性与局部单调性的概念。}

总之，若认为依从者中的效应值得关注，那么在双臂、全或无依从的双盲随机化试验中——尤其当其中一个臂按设计会表现出完全依从时——依赖单调性 (iv) 似乎是可行之道。然而，在更复杂的情形与观察性研究中使用这一方法时需要谨慎，即便所提出的工具真是工具。\mnote{Sommer 与 Zeger (1991)、Imbens 与 Rubin (1997) 以及 Greenland (2000) 描述了对照组完全依从的例子。}

\section{再论三个工具变量条件}

前几节讨论了选择单调性或同质性作为条件 (iv) 的相对优劣。我们的讨论隐含假定所提出的工具 $Z$ 事实上是工具。然而在观察性研究中，若所提出的工具 $Z$ 违背工具变量条件 (ii) 或 (iii) 中的任一条，它就不是有效工具；若它仅勉强满足条件 (i)，它就是弱工具。

在所有这些情形下，即便条件 (iv) 完全成立，使用 IV 估计也可能导致相当大的偏倚。下面我们逐一讨论三个工具变量条件。

条件 (i)，即 $Z$ 与 $A$ 的关联，可以凭经验检验。在宣布 $Z$ 为所提出的工具之前，研究者会检查 $Z$ 与处理 $A$ 是否相关。然而，当这一关联如我们的戒烟例子中那样微弱时，该工具被称为弱工具（见精要 16.2）。所提出的工具为弱工具时，会出现三个严重问题。

\begin{finepoint}[弱工具变量的定义]
文献中有两种相关但不同的弱工具定义：

\begin{enumerate}[leftmargin=1.9em,itemsep=0.35em,topsep=0.3em]
\item 若 $Z$ 与 $A$ 之关联的真值——IV 估计目标的分母——“很小”，则称工具（实质上）弱。
\item 若与观测到的 $Z$ 与 $A$ 之关联相对应的 $F$ 统计量“很小”，通常指小于 10，则称工具（统计上）弱。
\end{enumerate}

在我们的戒烟例子中，所提出的工具同时满足两种定义：风险差仅为 6\%，而 $F$ 统计量仅为 0.8。

第一种定义基于 $Z$ 与 $A$ 之关联的真值，它提醒我们：在使用弱候选工具时，即便样本量无穷大，IV 估计量也会极大地放大分子中的任何偏倚（即正文所述弱工具的第二问题）。第二种定义基于 $Z$ 与 $A$ 之关联的统计性质，它提醒我们：即便拥有完美的工具 $Z$，IV 估计量在有限样本中也可能有偏（即正文所述弱工具的第三问题）。
\end{finepoint}

第一，弱工具给出的效应估计置信区间很宽，如 16.2 节的戒烟例子所示。第二，弱工具会放大违背条件 (ii) 与 (iii) 所致的偏倚。与处理 $A$ 弱相关的候选工具 $Z$ 使 IV 估计量的分母很小。因此，影响 IV 估计量分子的、对条件 (ii) 与 (iii) 的违背（例如工具存在未测量混杂、工具存在直接效应）会被极大地夸大。在我们的例子中，任何影响 IV 估计量分子的偏倚都会被放大约 15.9 倍（$1/0.0627$）。第三，即便工具有效、样本量很大，弱工具也会在标准 IV 估计量中引入偏倚。

为理解第三个问题的性质，考虑一个随机生成的二值变量 $Z$。\mnote{在线性模型中，Martens 等 (2006) 表明：强混杂会使弱工具必然弱。}在无限总体中，IV 估计目标的分母将恰好为零——处理 $A$ 与一个完全随机的变量之间零关联——IV 估计值无定义。然而在有限样本的研究中，偶然性会使随机生成的 $Z$ 与未测量混杂因子 $U$ 之间、从而 $Z$ 与处理 $A$ 之间，出现微弱但不恰好为零的关联。若我们把这一随机 $Z$ 当作工具，IV 估计量的分母将非常小而非零。结果，分子被错误地放大，可能产生非常大的偏倚。事实上，我们所提出的工具“价格高于 1.50 美元”表现得就像一个随机生成的变量。\mnote{代码：程序 16.4}若我们把 $Z$ 定义为价格高于 1.60、1.70、1.80 或 1.90 美元，IV 估计值将分别为 $41.3$、$-40.9$、$-21.1$ 或 $-12.8$ kg。每种情形下，估计值周围的 95\% 置信区间都非常大。鉴于弱工具可能造成多大的偏倚与变异，一个略微违背条件 (ii) 与 (iii) 的强候选工具，可能优于一个违背较少但更弱的候选工具。\mnote{Bound、Jaeger 与 Baker (1995) 记录了这一偏倚。他们的论文之后，许多其他论文研究了弱工具的缺陷。}

条件 (ii)，即工具对结局没有直接效应，无法由数据验证。对条件 (ii) 的偏离，可以表示为从工具 $Z$ 指向结局 $Y$ 的一条直接箭头，如图 16.8 所示。这一不经处理 $A$ 中介的工具直接效应会进入 IV 估计量的分子，并被分母错误地放大，仿佛它是处理 $A$ 之效应的一部分。

\begin{figure}[!ht]
\centering
\begin{tikzpicture}[causal]
  \node[vtx] (Z) at (0,0) {$Z$};
  \node[vtx] (A) at (1.9,0) {$A$};
  \node[vtx] (Y) at (3.8,0) {$Y$};
  \node[vtx] (U) at (0.5,-1.5) {$U$};
  \draw[ed] (Z) -- (A);
  \draw[ed] (A) -- (Y);
  \draw[ed] (U) -- (A);
  \draw[ed] (U) -- (Y);
  \draw[ed] (Z.north) to[bend left=40] (Y.north);
\end{tikzpicture}
\caption{}
\label{fig:16.8}
\end{figure}

当分析中用较粗（例如二分）版本的 $A^{*}$ 替代连续或多值处理 $A$ 时，条件 (ii) 可能被违背。图 16.9 表明，即便条件 (ii) 对原始处理 $A$ 成立，它对二分版本 $A^{*}$ 也未必成立，因为路径 $Z\rightarrow A\rightarrow Y$ 表示工具 $Z$ 的一个直接效应，它不经 IV 分析所要估计其效应的那个处理 $A^{*}$ 中介。实践中，为解释简便，许多处理都被替换为较粗的版本。处理的粗化对 IV 估计是个问题，但对此前各章讨论的方法未必是问题。

\begin{figure}[!ht]
\centering
\begin{tikzpicture}[causal]
  \node[vtx] (Ast) at (1.9,1.5) {$A^{*}$};
  \node[vtx] (Z) at (0,0) {$Z$};
  \node[vtx] (A) at (1.9,0) {$A$};
  \node[vtx] (Y) at (3.8,0) {$Y$};
  \node[vtx] (U) at (0.5,-1.5) {$U$};
  \draw[ed] (Z) -- (A);
  \draw[ed] (A) -- (Ast);
  \draw[ed] (A) -- (Y);
  \draw[ed] (U) -- (A);
  \draw[ed] (U) -- (Y);
\end{tikzpicture}
\caption{}
\label{fig:16.9}
\end{figure}

条件 (iii)，即工具对结局之效应无混杂，同样无法检验。图 16.10 显示了由候选工具 $Z$ 与结局 $Y$ 的共同原因所致的混杂，这些共同原因可能是（$U_1$）也可能不是（$U_2$）处理 $A$ 的原因。在观察性研究中，候选工具存在混杂的可能性始终存在（与任何其他不受研究者控制的变量一样）。混杂会进入 IV 估计量的分子，并被分母错误地放大，仿佛它是处理 $A$ 对结局 $Y$ 之效应的一部分。

\begin{figure}[!ht]
\centering
\begin{tikzpicture}[causal]
  \node[vtx] (U2) at (0.5,1.5) {$U_2$};
  \node[vtx] (Z) at (0,0) {$Z$};
  \node[vtx] (A) at (1.9,0) {$A$};
  \node[vtx] (Y) at (3.8,0) {$Y$};
  \node[vtx] (U1) at (0.5,-1.5) {$U_1$};
  \draw[ed] (U2) -- (Z);
  \draw[ed] (U2) -- (Y);
  \draw[ed] (U1) -- (Z);
  \draw[ed] (U1) -- (A);
  \draw[ed] (U1) -- (Y);
  \draw[ed] (Z) -- (A);
  \draw[ed] (A) -- (Y);
\end{tikzpicture}
\caption{}
\label{fig:16.10}
\end{figure}

有时，条件 (iii) 以及其他条件，在可测量协变量的各层内可能显得更可信。与其作出“$Z$ 对 $Y$ 的效应绝对没有混杂”这一无法检验的假设，我们或许更愿意作出“在可测量的工具前协变量 $V$ 的各层内，$Z$ 对 $Y$ 的效应没有未测量混杂”这一无法检验的假设。于是我们可以在 $V$ 的每一层内重复应用 IV 估计，并在“总体中的效应（同质性下）或依从者中的效应（单调性下）在 $V$ 各层内恒定”的假设下合并各层的 IV 效应估计。另一种做法是把变量 $V$ 作为协变量纳入两阶段建模。在我们的例子中，这减小了效应估计的绝对值，并增大了其 95\% 置信区间。\mnote{代码：程序 16.5}

支持条件 (iii) 的另一种常见策略，是检查可测量混杂因子在候选工具 $Z$ 各层间的分布是否均衡。其思路是：若可测量的混杂因子均衡，则未测量的混杂因子也更可能均衡。然而，这种做法可能带来虚假的安全感：由于上述偏倚放大，即便很小的不均衡也可能导致反直觉的大偏倚。

即便 $Z$ 对 $Y$ 的效应没有混杂，条件 (iii) 也可能被违背。条件 (iii) 的正式版本要求取值不同的候选工具者之间可交换。这种可交换性可能因混杂（见上）或选择偏倚而被违背。IV 分析中引入选择偏倚的一种出人意料地常见的方式，是排除取某些处理 $A$ 值的个体。例如，若总体中的个体可能接受 $A=0$、$A=1$ 或 $A=2$ 的处理水平，那么仅限于 $A=1$ 或 $A=2$ 者的 IV 分析可能给出非零的效应估计，即便真实因果效应为零。对于以估计 $A=1$ 相对 $A=2$ 之效应为目标的非 IV 分析，这种排除并不会引入偏倚。\mnote{Swanson 等 (2015c) 详细描述了这一偏倚。}

上述与条件 (i)–(iii) 相关的全部问题，在同时使用多个候选工具以图缓解单一候选工具之弱性的 IV 分析中都会加剧。遗憾的是，候选工具的数量越多，其中某些工具违背某条工具变量条件的可能性就越大。

\section{工具变量估计与其他方法的比较}

IV 估计至少在三个方面不同于此前讨论的所有方法。

第一，IV 估计用其他假设取代了条件可交换性假设。IP 加权与标准化要求受治疗者与未受治疗者在给定可测量变量下可交换。相反，当条件 (i)–(iv) 成立时，即便条件可交换性不成立，IV 估计也能提供有效的效应估计。因此，方法的选择取决于：在特定的研究情境中，是识别并测量 $A$ 对 $Y$ 之效应的混杂因子更容易，还是找到工具 $Z$ 并指望单调性或不存在相关效应异质性更容易。

第二，对 IV 估计的条件 (i)–(iv) 相对轻微的违背，可能导致很大的偏倚。\mnote{IV 估计不是唯一不需要条件可交换性的方法：回归断点（精要 16.3）与双重差分（技术注记 7.3）亦然。}IV 估计的基础是分母放大分子。因此，当条件不完美成立或工具较弱时，两个方向上都有出现严重偏倚的可能。结果，IV 估计有时可能比未校正的估计偏倚更大。相反，当可识别性条件仅被轻微违背时，IP 加权与标准化往往只给出略有偏倚的估计，校正引入大偏倚的可能性较小。IV 估计对其可识别性条件之偏离的敏感性，凸显了敏感性分析的重要性。\mnote{Baiocchi 等 (2014) 综述了一些量化 IV 估计对关键假设之违背的敏感程度的方法。}

第三，标准 IV 估计之适用性的理想情境，比其他方法更为受限。如本章所论，标准 IV 估计最好留给这样的情境：存在大量未测量混杂；处理 $A$ 真正是二值且时点固定的；候选工具 $Z$ 强（且为因果工具）；并且预期同质性成立，或者——若确实关心依从者中的效应——单调性成立。这些限制的一个后果是，IV 估计一般用于回答关于点干预的因果问题。因此，在本书第三部分——讨论时变处理与长期维持的复杂处理策略之对比——IV 估计不会成为主要方法。

因果推断依赖于假设的透明性，以及来自依赖不同假设集的方法之结果的三角互证。因此，IV 估计是一种有吸引力的做法，因为它依赖的假设集与其他方法不同。然而，由于 IV 估计典型的宽 95\% 置信区间，使用这一方法带来的增值往往很小。此外，IV 估计的使用者需要批判性地意识到该方法的局限。虽然这句话显然适用于任何因果推断方法，但 IV 估计中偏倚在方向上与量级上可能反直觉，需要格外注意。\mnote{透明性要求对 IV 分析作恰当的报告。关于建议的报告准则，见 Brookhart 等 (2010)、Swanson 与 Hernán (2013) 以及 Baiocchi 等 (2014)。}

\begin{techpoint}[更一般的结构均值模型]
考虑一个允许连续和（或）多值处理 $A$、工具 $Z$ 与工具前协变量 $V$ 的可加结构均值模型。这类模型假定
\[
\E[Y-Y^{a=0}|Z,A,V]=\gamma(Z,A,V;\beta),
\]
其中 $\gamma(Z,A,V;\beta)$ 是已知函数，$\beta$ 是未知的（可能为向量值的）参数，且 $\gamma(Z,A=0,V;\beta)=0$。也就是说，可加结构均值模型是关于“处理水平 $A$ 相对处理水平 0 之平均因果效应”的模型，其对象是那些工具取 $Z$、混杂因子取 $V$、且观测到的处理恰为 $A$ 的个体子集。在条件反事实均值独立性假设 $\E[Y^{a=0}|Z=1,V]=\E[Y^{a=0}|Z=0,V]$ 下，该模型的参数可用 g 估计识别。

类似地，一般的可乘结构均值模型假定
\[
\E[Y|Z,A,V]=\E[Y^{a=0}|Z,A,V]\exp[\gamma(Z,A,V;\beta)],
\]
其中 $\gamma(Z,A,V;\beta)$ 是已知函数，$\beta$ 是未知参数向量，且 $\gamma(Z,A=0,V;\beta)=0$。在类似条件下，该模型的参数也可用 g 估计识别。结构均值模型的识别条件与有效估计量由 Robins (1994) 讨论，并由 Vansteelandt 与 Goetghebeur (2003) 综述。更一般地，嵌套可加与可乘结构均值模型的 g 估计，可以把针对时点固定处理与混杂因子的 IV 方法推广到时变处理与混杂因子的情形。
\end{techpoint}

\begin{techpoint}[单调性与依从者中的效应]
考虑一个二值因果工具 $Z$（如正文所述的随机化指示变量）与处理 $A$。Imbens 与 Angrist (1994) 证明：在单调性 (iv) 下，即不存在违抗者时，通常的 IV 估计目标等于依从者中的平均因果效应 $\E[Y^{a=1}-Y^{a=0}|A^{z=1}-A^{z=0}=1]$。Baker 与 Lindeman (1994) 对二值结局给出了相关证明。另见 Angrist、Imbens 与 Rubin (1996) 及随后的讨论，以及 Baker、Kramer 与 Lindeman (2016)。证明如下。

处理分派的效应（意向性治疗效应）可以写作四个主层中意向性治疗效应的加权平均，式中四项依次对应永远接受者、永不接受者、依从者与违抗者：
\begin{footnotesize}
\begin{align*}
\E[Y^{z=1}-Y^{z=0}] ={}& \E\bigl[Y^{z=1}-Y^{z=0}\bigm|A^{z=1}=1,A^{z=0}=1\bigr]\Pr[A^{z=1}=1,A^{z=0}=1]\\
&+\E\bigl[Y^{z=1}-Y^{z=0}\bigm|A^{z=1}=0,A^{z=0}=0\bigr]\Pr[A^{z=1}=0,A^{z=0}=0]\\
&+\E\bigl[Y^{z=1}-Y^{z=0}\bigm|A^{z=1}=1,A^{z=0}=0\bigr]\Pr[A^{z=1}=1,A^{z=0}=0]\\
&+\E\bigl[Y^{z=1}-Y^{z=0}\bigm|A^{z=1}=0,A^{z=0}=1\bigr]\Pr[A^{z=1}=0,A^{z=0}=1].
\end{align*}
\end{footnotesize}

然而，永远接受者与永不接受者中的意向性治疗效应都为零，因为 $Z$ 在这两个主层中不影响 $A$，且由技术注记 16.1 的个体层面条件 (ii)，$Z$ 对 $Y$ 没有独立效应。若假定不存在违抗者，则上述和式简化为依从者一项：
\[
\E[Y^{z=1}-Y^{z=0}]=\E\bigl[Y^{z=1}-Y^{z=0}\bigm|A^{z=1}=1,A^{z=0}=0\bigr]\Pr[A^{z=1}=1,A^{z=0}=0].
\]
但在依从者中，$Z$ 对 $Y$ 的效应等于 $A$ 对 $Y$ 的效应（因为 $Z=A$），即
\begin{align*}
&\E\bigl[Y^{z=1}-Y^{z=0}\bigm|A^{z=1}=1,A^{z=0}=0\bigr]\\
&\qquad=\E\bigl[Y^{a=1}-Y^{a=0}\bigm|A^{z=1}=1,A^{z=0}=0\bigr].
\end{align*}
因此，依从者中的效应为
\[
\frac{\E[Y^{z=1}-Y^{z=0}]}{\Pr[A^{z=1}=1,A^{z=0}=0]}=\E\bigl[Y^{a=1}-Y^{a=0}\bigm|A^{z=1}=1,A^{z=0}=0\bigr],
\]
若假定 $Z$ 是随机分配的，它就是通常的 IV 估计目标；因为随机分配蕴含 $Z\indep\{Y^{a,z},A^{z};z=0,1;a=0,1\}$。在这一联合独立性与一致性下，分子中的意向性治疗效应等于 $\E[Y|Z=1]-\E[Y|Z=0]$，分母中依从者的比例等于 $\Pr[A=1|Z=1]-\Pr[A=1|Z=0]$。要看后一等式为何成立，注意永远接受者的比例 $\Pr[A^{z=1}=1,A^{z=0}=1]=\Pr[A=1|Z=0]$，永不接受者的比例 $\Pr[A^{z=1}=0,A^{z=0}=0]=\Pr[A=0|Z=1]$。由于在单调性 (iv) 下不存在违抗者，依从者的比例等于余下的部分
\begin{align*}
1-\Pr[A=1|Z=0]&-\Pr[A=0|Z=1]\\
&=1-\Pr[A=1|Z=0]-(1-\Pr[A=1|Z=1])\\
&=\Pr[A=1|Z=1]-\Pr[A=1|Z=0],
\end{align*}
这就完成了证明。

上述证明只考虑了图 16.1 所描绘的、工具 $Z$ 为因果工具的情形。当如图 16.2 与 16.3 所示，我们只有替代工具 $Z$ 的数据、而没有因果工具 $U_Z$ 的数据时，Hernán 与 Robins (2006b) 证明：依从者（按 $U_Z$ 定义）中的平均因果效应同样可由通常的 IV 估计量识别。他们的证明关键依赖于两个假设：给定因果工具 $U_Z$ 时 $Z$ 独立于 $A$ 与 $Y$，以及 $U_Z$ 是二值的。然而，除非 $U_Z$ 是连续的，这一独立性假设往往缺乏实质可信性。一个推论是：在许多把 IV 方法用于观察性数据的应用中，$Z$ 充其量只是 $U_Z$ 的替代，此时把 IV 估计目标解释为依从者中的效应值得怀疑。
\end{techpoint}

\begin{finepoint}[回归断点设计]
假设我们关心一种新的抗病毒治疗 $A$ 对 1 周后测量的连续结局——血氧水平 $Y$ 的效应。凡是因诊断为 COVID-19 到达医院的人，都适用该治疗。然而，由于治疗供应短缺，卫生当局禁止对 65 岁以下者施用该治疗，以保证 65 岁及以上者都能得到治疗。也就是说，接受治疗的概率为 $\Pr[A=1|L<65]=0$、$\Pr[A=1|L\ge 65]=1$，其中 $L$ 是年龄。这里没有正值性：受治疗者与未受治疗者没有重叠的混杂因子 $L$ 取值。

在缺乏正值性时，我们需要作出替代假设来识别因果效应。一个合理的假设是：给定 $L$ 时反事实结局的条件均值 $\E[Y^{a=1}|L]$ 与 $\E[Y^{a=0}|L]$ 关于 $L$ 连续。换言之，若我们能沿年龄轴画出这两个均值（我们做不到，因为这些均值是反事实的，因而未被观测），我们不会在线条上看到任何跳跃。在这一连续性假设，以及“阈值两侧附近的个体可比较”这一可交换性假设之下，给定 $L$ 时观测结局之条件均值 $\E[Y|L]$ 在 $L=65$ 附近的断点，就可以解释为治疗概率在 65 岁处突变的结果。$Y$ 的均值是否在阈值处跳跃，可以通过画出观测数据凭经验检查。（严格说来，我们只需在阈值 $L=65$ 附近连续即可。）

因此，在连续性假设下，我们可以把 $A$ 的平均因果效应估计为：阈值上方附近个体（比如 65 岁零 1 个月者）的平均结局 $Y$ 与阈值下方附近个体（比如 64 岁零 11 个月者）的平均结局之差。若阈值周围 1 个月的带宽太小（因为数据中处于该范围的个体太少），我们就需要增大阈值周围的带宽。例如，我们可以用 1 年的带宽，比较 64 岁者与 65 岁者。带宽的选择至关重要：较宽的年龄区间可能因比较不可交换的个体而引入偏倚。带宽一旦确定，我们在阈值 $L=65$ 两侧分别拟合线性回归模型，以估计阈值两侧的平均结局。为帮助确定阈值周围的带宽，可以使用交叉验证之类数据自适应程序（见精要 18.2）。此外，若认为达到条件可交换性有必要，回归模型中可以纳入协变量。

上述方法称为\emph{回归断点设计}（regression discontinuity design），由 Thistlewaite 与 Campbell (1960) 首先提出。当用单个协变量 $L$ 来分派治疗，且在“$L$ 与 $Y^{a}$ 的关系平滑（即无跳跃）”这一连续性假设下，它就可以使用。回归断点设计估计的是总体中 $L$ 取值接近阈值的那个子集中，处理 $A$ 对结局 $Y$ 的平均因果效应。若 $L$ 是效应修饰因子，这一条件效应可能不同于总体中的平均因果效应。注意，若除 $A$ 之外的其他治疗也在阈值附近发生变化（例如卫生当局同样把稀缺的重症监护床位限制给 65 岁及以上者），或者知道阈值的高风险个体会操纵自己的数据（例如 63、64 岁的高风险者想办法提供显示更老年龄的假文件），回归断点设计将给出条件效应的有偏估计。

更具体地说，我们这里描述的是\emph{锐回归断点设计}（sharp regression discontinuity design），其中治疗概率在阈值处从 0 跳到 1。\emph{模糊回归断点设计}（fuzzy regression discontinuity design）是这一方法的扩展，它允许治疗概率从大于 0 的某个值跳到小于 1 的某个值。这一扩展依赖单调性假设，估计的是总体中一个子集的子集——即 $L$ 取值接近阈值的依从者——中的平均因果效应。关于估计的细节，见 Hahn、Todd 与 van der Klaauw (2001) 以及 Imbens 与 Lemieux (2008)。
\end{finepoint}
